\chapter{Матричный обменный мост и знак ориентационного взаимодействия}

Предыдущая проверка установила, что для рождения $\mathbb{RP}^2$ необходим
отрицательный коллективный инвариант чистоты. Теперь скалярная амплитуда
обменного моста заменяется общей вещественной семейной матрицей, чтобы
проверить, не скрыт ли этот инвариант в уже выведенном каноническом
действии.

\section{Единственное естественное матричное продолжение}

Коэффициентный проектор ранга пятнадцать раскладывается как произведение
пятимерной кратности и семейного триплета. Поэтому скаляр
$\lambda I_3$ можно заменить на $B\in M_3(\mathbb R)$. Два дефекта
унитарности становятся $I_3-B^TB$ и $I_3-BB^T$. Нормировка, точно
воспроизводящая скалярное действие, равна
\begin{equation}
 \mathcal S_{\rm br}(B)
 =\frac1{21}\left[
  \Tr(I_3-B^TB)^2+\Tr(I_3-BB^T)^2\right]
 =\frac2{21}\Tr(I_3-B^TB)^2.
 \label{eq:v6-matrix-bridge-action}
\end{equation}
При $B=\lambda I_3$ получаем прежнее
$2(1-\lambda^2)^2/7$ без нового коэффициента.

Пусть $X=B^TB\ge0$. Тогда
\begin{equation}
 \mathcal S_{\rm br}
 =\frac2{21}\left(3-2\Tr X+\Tr X^2\right).
 \label{eq:v6-bridge-purity-sign}
\end{equation}
При фиксированном $r=\Tr X$ неравенство Коши даёт
\begin{equation}
 \Tr X^2\ge\frac{r^2}{3},
\end{equation}
причём равенство достигается только при $X=(r/3)I_3$.

\begin{theorem}[Канонический мост подавляет одноосность]
Матричное продолжение канонического положительного действия минимизирует
чистоту $X$ при фиксированной общей амплитуде и выбирает три равных
сингулярных значения. Следовательно, оно стабилизирует изотропную фазу,
а не индуцирует расщепление $3\to1+2$.
\end{theorem}

\section{Точный словарь к состоянию}

На слое $\Tr X=3$ естественная нормировка состояния есть
\begin{equation}
 R=\frac13X.
\end{equation}
Тогда \eqref{eq:v6-bridge-purity-sign} превращается в
\begin{equation}
 \boxed{
 \mathcal S_{\rm br}(R)
 =\frac67\left(\Tr R^2-\frac13\right)}.
 \label{eq:v6-bridge-positive-purity}
\end{equation}
Это ровно инвариант предыдущей проверки, но с противоположным знаком.
В соглашении
$\Tr(R\log R)-\kappa(\Tr R^2-1/3)$ мост соответствует
\begin{equation}
 \kappa_{\rm br}=-\frac67,
\end{equation}
а не положительному $\kappa>\log4$.

Для критического одноосного состояния
\begin{equation}
 R_c=\frac23P+\frac16(I_3-P)
\end{equation}
получается точная цена
\begin{equation}
 \mathcal S_{\rm br}(R_c)
 =\frac67\left(\frac12-\frac13\right)
 =\frac17.
 \label{eq:v6-critical-state-bridge-cost}
\end{equation}
Совпадение с топологическим весом $1/7$ точно, но его смысл ограничен:
оно использует условную идентификацию $R=X/3$ и является стоимостью
отклонения от замкнутого моста, а не выигрышем материальной фазы.

\section{Порог для возможной квантовой поправки}

Если будущий флуктуационный член породит
$-\kappa_{\rm fluc}(\Tr R^2-1/3)$ в той же нормировке, то эффективный
коэффициент равен
\begin{equation}
 \kappa_{\rm eff}=\kappa_{\rm fluc}-\frac67.
\end{equation}
Для перехода требуется
\begin{equation}
 \kappa_{\rm fluc}>\log4+\frac67
 =2.243437218\ldots.
 \label{eq:v6-fluctuation-alignment-threshold}
\end{equation}
Это не предсказание коэффициента, а строгий kill-threshold для следующего
кандидата.

\begin{nogocriterion}[Запрет смены знака моста]
Нельзя использовать матричный обменный мост как источник ориентационной
фазы, меняя знак его положительной нормы. Допустим только независимо
выведенный флуктуационный или аномальный член, который в одной и той же
нормировке преодолевает порог
$\log4+6/7$.
\end{nogocriterion}

\begin{protocol}
\begin{enumerate}
 \item матричное продолжение скалярного моста:
 \textbf{построено без нового коэффициента};
 \item восстановление скалярной формулы:
 \textbf{точное};
 \item знак инварианта $\Tr R^2$:
 \textbf{положительный};
 \item предпочтение при фиксированном следе:
 \textbf{изотропное состояние};
 \item индуцирование $-\Tr R^2$:
 \textbf{нет};
 \item цена критического одноосного состояния:
 \textbf{$1/7$};
 \item достаточность этого совпадения для рождения материи:
 \textbf{нет};
 \item минимальный требуемый флуктуационный коэффициент:
 \textbf{$\log4+6/7$ в данной нормировке};
 \item рождение $\mathbb{RP}^2$ текущим мостом:
 \textbf{закрыто отрицательно};
 \item следующая проверка:
 \textbf{флуктуационный детерминант матричного моста и знак чистоты}.
\end{enumerate}
\end{protocol}

Машинный сертификат сохранён в
\path{s2t_v6_exchange_bridge_induced_alignment_gate_results.json}.